\begin{proposition}[P4.15]\hypertarget{tag-02WM}{}\label{night:P4.15}\label{stmt:P4.15}
\StatementRecord{night:P4.15}{proved}{02WM}{\SeriesRef{night:P4.12}, \SeriesRef{night:P4.13}, \SeriesRef{night:P4.8}}

Let $mS$ be an irreducible singular fibre in the setting of 4.2.
Suppose its normalization $T$ is smooth and $d\nu:T_T\to\nu^*T_X$
has rank two away from finitely many points. Let $\ell$ be the
length of its maximal-minor scheme. Then
\[
 e(T)=K_T^2-\ell.
\]
The alternatives for $T$ are a rational surface with
$e(T)=6-\ell/2$ and $0\leq\ell\leq6$, a minimal ruled surface
over an elliptic curve with $\ell=0$, or a Hopf surface with
$\ell=0$. In either of the last two alternatives $e(S)=0$.
The reduced conductor and singular locus are disjoint unions of
smooth elliptic curves, and their normalization maps are finite
and unramified.
\end{proposition}
\begin{proof}

\begin{enumerate}[label={$\langle1\rangle$\arabic*.},leftmargin=*]
\item The differential is injective as a sheaf map: its kernel vanishes
on a dense open set, and a locally free sheaf is torsion-free.
Its rank-one cokernel $Q$ is locally free away from finitely many
points. The depth lemma applied to the length-one locally free
resolution gives depth at least one at those points, excluding
a finite-length subsheaf. All torsion would be supported there,
so $Q$ is torsion-free. On the smooth surface we obtain
\[
 Q=\mathcal I_Z\otimes L,\qquad
 L=K_T\otimes\nu^*(\det T_X),\qquad \operatorname{length}(Z)=\ell.
\]
The two-column presentation by $d\nu$ identifies the local ideal
of $Z$ with its three maximal minors. The classes $c_1(T_X)$ and
$c_2(T_X)$ vanish in real cohomology, since $b_2(X)=b_4(X)=0$.
The second Chern class of the exact sequence therefore gives
\[
 0=e(T)-K_T^2+\ell.
\]
P4.12 gives $\kappa(T)=-\infty$. For a ruled surface over a
genus-$g$ curve, blown up $k$ times, we have
\[
 e(T)=4(1-g)+k,\qquad K_T^2=8(1-g)-k.
\]
Consequently $4(g-1)+2k+\ell=0$. Thus $g\leq1$, and $g=1$
forces $k=\ell=0$. In the class-VII case, including blowups,
$K_T^2=-e(T)$ and $e(T)=b_2(T)\geq0$. Hence
$2e(T)+\ell=0$, and the surface is minimal with $b_2=0$.
Its nonempty conductor, the finite-cover classification, and the
absence of curves on Inoue surfaces proved in P4.8 force it to
be Hopf. For rational $T$, Noether's formula gives
$2e(T)+\ell=12$. The inequality $e(T)\geq3$ gives the asserted
range. When $\ell=0$, the cokernel is a line bundle and $d\nu$
has rank two everywhere.

\item We now compute the image Euler number in the two latter cases.
The conductor divisor $D$ satisfies $K_T+D\equiv0$. First let
$T$ be a minimal ruled surface over an elliptic curve $E$;
it is a $\mathbb P^1$-bundle. For a ruling fibre $F$, $D\cdot F=2$.
If a horizontal component $H$ occurs with multiplicity one,
adjunction gives
\[
 2p_a(H)-2=(K_T+H)H=-(D-H)H\leq0.
\]
The map to $E$ gives geometric genus at least one. Equality
therefore holds, $H$ is smooth elliptic, and $H$ is disjoint
from every other component of $D$. There is no vertical
component because every ruling fibre meets $H$. The possibilities
are one elliptic double-section or two disjoint elliptic sections.

\item The other possibility is $D=2H+\sum_j a_jF_j$, where $H$ is a
section, $F_j$ are distinct ruling fibres and $a_j>0$.
The reduced conductor has Euler number equal to the number $k$
of these fibres: $e(H)=0$, $e(F_j)=2$, and each meets $H$ once.
The immersive normalization identifies $T\setminus D_{\mathrm{red}}$
with $S_{\mathrm{reg}}$. Hence
\[
 e(S_{\mathrm{reg}})=e(T)-e(D_{\mathrm{red}})=-k.
\]
P4.13 forces $k=0$. Thus in every case the reduced conductor is
a disjoint union of smooth elliptic curves.

\item On a Hopf surface every irreducible curve $H$ satisfies
$H^2=K_TH=0$, so $p_a(H)=1$. A rational curve would lift from
$\mathbb P^1$ to the universal cover $\mathbb C^2\setminus\{0\}$,
where both coordinate functions must be constant. Thus $H$
is smooth elliptic. Distinct curves have intersection zero and
are disjoint, proving the same conductor description.

\item To check the image, at a singular point the finite immersive
normalization presents the germ as distinct smooth surface
branches $S_i$. The reduced conductor on each branch is a single
smooth curve germ $C_i$. Any two branches intersect in a curve
in the smooth ambient threefold, and its support lies in both
$C_i$ and $C_j$. Their irreducibility forces those curve germs
to agree. Thus the image singular locus is one smooth curve
germ, and the conductor maps locally isomorphically to it.
Globally this map is finite and unramified. Riemann--Hurwitz
then makes every image component elliptic. Finally
\[
 e(S)=e(T)-e(D_{\mathrm{red}})+e(\operatorname{Sing}S)=0.
\]

\end{enumerate}
\end{proof}